\section{An Operational Traceability Contract}
\subsection{Why an apparent vertical chain needs two additional tests}
The intuitive phrase ``from strategy to university rules'' may be read in three ways. The weakest, investigated here, is that suitable public instruments can be found at several governance levels. A stronger documentary statement is that an identifiable lower-level instrument \emph{expressly draws authority or control requirements} from a named upper-level one. Stronger still, actual teaching, assessment or data-handling behaviour might conform to those rules. These are distinct measurement targets, even when their narrative descriptions are similar.

For each jurisdiction and layer, let $A_{\ell j}$ denote the proposition that an authority- and scope-qualified public instrument is located. Then $L_{\ell j}=\stE$ is a recorded positive witness for $A_{\ell j}$ under Table~\ref{tab:gates}; $\stP$ and $\stU$ encode incomplete \emph{evidence}, not the falsity of the proposition. For an adjacent pair of layers define a separate documentary-link proposition $R_{\ell,\ell+1,j}$, requiring an identifiable source-to-target relationship, e.g., a lower instrument that explicitly invokes the upper instrument, a dated delegation, or a controlling sector circular with named institutional applicability. Define an implementation proposition $I_{\ell j}$ requiring records of use, compliance or enforcement in the relevant institutional setting. Our reported chain outcomes concern $A$; the $R$ and $I$ variables were not systematically examined or coded. They are \emph{unassessed}, not coded ``no'' and not automatically inferred from $A$.

This separation matters even for Saudi Arabia's fully explicit four-layer vector. Four admissible instruments do not demonstrate three directed links. Conversely, an L3 source unlocated in this public search does not rule out a substantive implementation pathway through unpublished ministerial communications. A later study could extend, rather than overwrite, the present frame by building an evidence graph with instrument-version nodes and directed reference edges. Each claimed edge should carry both the source and target ID, a quoted authority/obligation passage, the exact language/edition and the date when the edge was operative. Without this information, a schematic arrow is illustrative rather than an observed transmission fact.

\begin{table*}[t]
\centering\caption{Evidence contract separating available instruments, explicit documentary links, and implementation. Only the first object is measured in the present GCC results.}\label{tab:contract}
\footnotesize
\begin{tabularx}{\textwidth}{@{}p{0.16\textwidth}p{0.29\textwidth}p{0.29\textwidth}X@{}}
\toprule
Object & Minimum positive witness & What the witness supports & What it cannot independently establish\tabularnewline\midrule
Layer instrument $A_{\ell j}$ & Official, dated, inspectable instrument meeting the layer's authority and scope gate & \emph{Public existence} at that layer, coded $\stE/\stP/\stU$ & Explicit transmission to an adjacent document or any downstream behavior\tabularnewline
Directed link $R_{\ell,\ell+1,j}$ & Instrument-to-instrument citation, directive, delegation or explicit application clause & A specific \emph{documentary reference edge} between named instrument versions & Effective compliance or its substantive impact; \emph{not coded here}\tabularnewline
Implementation $I_{\ell j}$ & Auditable enforcement records, observed assessment practice, records review, or field interviews & Some \emph{realized rule operation} in a defined setting/time period & Nation-wide causal adoption; \emph{not observed here}\tabularnewline
\bottomrule
\end{tabularx}
\end{table*}

\subsection{Admission is not equivalent across layers}
A comparison is admissible only after naming its estimand. At L3, the question is whether a sector-applicable governance instrument exists with content that actually controls AI use in higher education. At L4, the question is whether \emph{any one} institution has issued such a rule for its own users. The former is a jurisdictional \emph{HE-sector instrument}; the latter a jurisdictional \emph{existence witness}. Replacing $A_{3j}$ with the count of local university documents would make the original gap disappear by definition while answering a different research question. Likewise, insisting that all universities must publish a rule to earn L4 $\stE$ would create a prevalence gate for which the study never assembled a representative denominator. The asymmetric gate is thus purposeful but must be stated at every interpretation of the L3--L4 contrast.

The approach is designed to be conservative about instrument authority while permissive about discovery. A direct official policy in one language is retained even if there is no translation; a second-language mirror cannot count as a second case. Where an official news article reports a proposed policy, it can be admitted as process evidence $\stP$ but does not receive the status of a promulgated control document. A course syllabus with a ChatGPT prohibition may evidence a binding \emph{course} rule, yet is not automatically an institution-wide GenAI governance instrument. The policy question becomes more precise, not merely more demanding: the coder records exactly \emph{which scope} the located document governs.
